\chapter{Metodologia}
\label{chap:metodologia}

Este capítulo apresenta a metodologia empírica adotada para investigar a
existência e as propriedades do prêmio de risco de volatilidade no mercado de
Bitcoin, bem como seu conteúdo informacional para a previsão de retornos futuros
do ativo. Com base nas evidências descritivas apresentadas no
Capítulo~\ref{chap:dados}, são formuladas as hipóteses de pesquisa e
especificados os modelos econométricos e preditivos utilizados ao longo da
dissertação.

A estratégia empírica combina modelos econométricos lineares de referência com
abordagens mais flexíveis, incluindo técnicas de \textit{machine learning}, com
o objetivo de capturar possíveis não linearidades, variação temporal de
parâmetros e interações complexas características do mercado de criptoativos.
Este capítulo concentra-se na definição formal das hipóteses, das variáveis e
das especificações empíricas, enquanto os resultados são apresentados nos
capítulos subsequentes.

% =========================================================
\section{Questões de pesquisa e hipóteses}
\label{sec:hipoteses}

A análise empírica desenvolvida nesta dissertação é guiada por duas questões
centrais. A primeira refere-se à existência de um prêmio de risco de
volatilidade estatisticamente e economicamente relevante no mercado de opções
de Bitcoin. A segunda diz respeito ao potencial conteúdo informacional desse
prêmio para a previsão de retornos futuros do Bitcoin no mercado à vista.

Com base na literatura sobre prêmios de risco de volatilidade em mercados
financeiros tradicionais e nas evidências descritivas apresentadas no
Capítulo~\ref{chap:dados}, são formuladas as seguintes hipóteses de pesquisa:

\begin{itemize}
    \item \textbf{H1:} O prêmio de risco de volatilidade do Bitcoin é, em média,
    estatisticamente diferente de zero ao longo da amostra, refletindo a
    compensação exigida pelos investidores para assumir risco de volatilidade
    no mercado de opções.

    \item \textbf{H2:} O prêmio de risco de volatilidade do Bitcoin contém
    informação relevante para a previsão de retornos futuros do ativo,
    apresentando relação estatisticamente significativa com os retornos
    esperados em horizontes subsequentes.
\end{itemize}

Essas hipóteses são avaliadas por meio de modelos econométricos lineares de
referência e de abordagens preditivas mais flexíveis, permitindo investigar
tanto a significância estatística quanto a relevância econômica do prêmio de
risco de volatilidade.

% =========================================================
\section{Definição das variáveis e horizontes de previsão}
\label{sec:variaveis_horizontes}

Nesta seção são definidas formalmente as variáveis utilizadas nos modelos
empíricos, bem como os horizontes temporais considerados na análise. Todas as
variáveis são construídas a partir das séries descritas no
Capítulo~\ref{chap:dados}, respeitando estritamente a informação disponível no
tempo $t$, de modo a evitar viés de \textit{look-ahead}.

A variável dependente principal é o retorno futuro do Bitcoin no mercado
\textit{spot}, definido como o retorno logarítmico acumulado ao longo de um
horizonte de $h$ dias:
\[
R_{t+h} = \sum_{i=1}^{h} r_{t+i},
\]
em que $r_t$ denota o retorno logarítmico diário e
$h \in \{1, 5, 10, 20, 30, 60\}$ representa o horizonte de previsão.

A principal variável explicativa é o prêmio de risco de volatilidade observado
no tempo $t$, definido como a diferença entre a volatilidade implícita extraída
do mercado de opções e a volatilidade realizada no mercado à vista, ambas
construídas em janelas móveis de 30 dias. Essa variável sintetiza a percepção de
risco e incerteza futura implícita nos preços das opções de Bitcoin.

Em especificações alternativas, são consideradas variáveis de controle
adicionais, incluindo medidas de volatilidade implícita e volatilidade
realizada isoladamente. Todas as variáveis explicativas são defasadas de forma
consistente, garantindo que estariam disponíveis aos agentes econômicos no
momento da formação das expectativas.

A consideração de múltiplos horizontes de previsão permite investigar se o
conteúdo informacional do prêmio de risco de volatilidade se manifesta de forma
distinta ao longo do tempo.

% =========================================================
\section{Modelos econométricos de referência}
\label{sec:modelos_referencia}

Esta seção apresenta os modelos econométricos de referência utilizados para
avaliar a relação entre o prêmio de risco de volatilidade do Bitcoin e os
retornos futuros do ativo. O objetivo desses modelos é estabelecer uma linha de
base empírica simples, transparente e comparável à literatura existente, antes
da introdução de abordagens mais flexíveis e não lineares.

Os modelos considerados são especificações lineares estimadas por mínimos
quadrados ordinários (OLS), nas quais os retornos futuros do Bitcoin são
explicados pelo prêmio de risco de volatilidade observado no período anterior.

\subsection{Especificação básica}

A especificação econométrica básica é dada por:
\[
R_{t+h} = \alpha_h + \beta_h \, \text{BVRP}_t + \varepsilon_{t+h},
\]
em que $R_{t+h}$ representa o retorno acumulado do Bitcoin no horizonte $h$,
$\text{BVRP}_t$ denota o prêmio de risco de volatilidade observado no tempo $t$,
$\alpha_h$ é um intercepto específico do horizonte de previsão,
$\beta_h$ é o coeficiente de interesse e $\varepsilon_{t+h}$ representa o termo
de erro.

O coeficiente $\beta_h$ captura a relação marginal entre o prêmio de risco de
volatilidade e os retornos futuros do Bitcoin. A análise empírica não impõe
restrições \textit{a priori} sobre o sinal desse coeficiente, permitindo que os
dados determinem a direção da relação.

As regressões são estimadas separadamente para cada horizonte de previsão,
permitindo avaliar como o efeito do prêmio de risco de volatilidade varia ao
longo do tempo.

\subsection{Extensão com variáveis de controle}

Para avaliar a robustez dos resultados e mitigar potenciais problemas de
omissão de variáveis relevantes, são consideradas extensões da especificação
básica que incluem variáveis de controle adicionais:
\[
R_{t+h} = \alpha_h + \beta_{1,h} \, \text{BVRP}_t + \beta_{2,h} \, \text{RV}_t
+ \beta_{3,h} \, \text{IV}_t + \varepsilon_{t+h},
\]
em que $\text{RV}_t$ e $\text{IV}_t$ representam, respectivamente, a
volatilidade realizada e a volatilidade implícita observadas no tempo $t$.

A inclusão dessas variáveis permite investigar se o prêmio de risco de
volatilidade possui poder explicativo adicional além do contido em medidas
tradicionais de volatilidade, reconhecendo a elevada correlação entre essas
variáveis.

\subsection{Inferência estatística}

A inferência estatística é conduzida com erros padrão robustos à
heterocedasticidade e à autocorrelação serial, reconhecendo que retornos
financeiros frequentemente exibem tais características. Em particular, são
utilizados erros padrão do tipo Newey--West, com número de defasagens compatível
com o horizonte de previsão considerado.

A significância estatística dos coeficientes é avaliada por meio de testes $t$,
enquanto a relevância econômica é analisada a partir da magnitude, sinal e
estabilidade dos coeficientes estimados ao longo dos diferentes horizontes e
especificações.

% =========================================================
\section{Estratégia empírica e procedimento de estimação}
\label{sec:estrategia_empirica}

A análise empírica é conduzida em frequência diária, respeitando a ordenação
temporal das informações. Em todas as especificações, as variáveis explicativas
são observadas no tempo $t$, enquanto as variáveis dependentes correspondem a
retornos acumulados em horizontes futuros $t+h$, assegurando que apenas
informação disponível \textit{ex ante} seja utilizada na previsão.

\subsection{Estrutura temporal e janelas de estimação}

Os modelos econométricos são estimados utilizando janelas móveis ao longo do
tempo, permitindo que os parâmetros se ajustem gradualmente a mudanças
estruturais no mercado de criptoativos. Esse procedimento reflete de forma mais
realista o ambiente de decisão enfrentado por agentes econômicos.

\subsection{Separação entre amostras de treino e teste}

Para as análises preditivas, a amostra é dividida em períodos de treino e teste,
respeitando estritamente a ordem temporal das observações. A amostra de treino é
utilizada para a estimação dos parâmetros, enquanto a amostra de teste é
reservada para a avaliação de desempenho fora da amostra.

\subsection{Sobreposição temporal}

Como os retornos futuros considerados correspondem a horizontes superiores a um
dia, as observações apresentam sobreposição temporal. Para lidar com esse
problema, a inferência estatística é conduzida utilizando erros padrão robustos
à autocorrelação, conforme descrito anteriormente.

% =========================================================
\section{Considerações sobre robustez e extensões}
\label{sec:robustez}

Além das especificações econométricas básicas, esta dissertação considera
análises de robustez destinadas a avaliar a estabilidade dos resultados em
diferentes horizontes de previsão, períodos amostrais e conjuntos de variáveis
explicativas.

Adicionalmente, são exploradas extensões que incorporam não linearidades,
interações entre variáveis e possíveis efeitos de regime, reconhecendo que a
relação entre volatilidade implícita, volatilidade realizada e retornos futuros
do Bitcoin pode variar ao longo do tempo. Essas extensões motivam a adoção de
modelos de \textit{machine learning}, desenvolvidos em detalhe no
Capítulo~\ref{chap:ml}.
